\documentclass[11pt]{article}
\usepackage[margin=1in]{geometry}
\usepackage{graphicx,amsmath,amssymb,booktabs,hyperref,xcolor}
\graphicspath{{/tmp/cheby_h0/figs/}}
\hypersetup{colorlinks=true,linkcolor=blue,urlcolor=blue}
\title{Methods and results: toward machine-$\varepsilon$ strict-$h{=}0$\\
       K=3 CRRA REE on the Chebyshev tabulated architecture}
\author{Fixed-Point Factory --- auto-mode running notes}
\date{June 2026}

\begin{document}
\maketitle

\begin{abstract}
\noindent
Cumulative document of methods explored to drive the K=3 CRRA REE
fixed-point residual to machine $\varepsilon$ on the tensor-Chebyshev
tabulated-$\mu$ architecture. Two clean wins so far:
\textbf{(1) kernel-band Cheb-tab nails to $\sim\!10^{-16}$ from cold and
warm} at bandwidths $h \in [0.20, 0.50]$; the FP depends on $h$ and the
joint $(N, h) \to (\infty, 0)$ limit gives slope $\alpha^* \approx 0.22$;
\textbf{(2) partition-of-unity (POU) Cheb-tab at strict $h{=}0$
recovers the correct REE slope $\alpha^* \approx 0.37$ at $N{=}8$},
matching the legacy \texttt{hfree\_smooth} at G=9. The POU operator
floors at $\sim 4\!\times\!10^{-3}$ with generic Anderson+Newton-Krylov;
several routes to lower (chebroots all-roots, analytic/complex-step
Jacobian) are under investigation.
\end{abstract}

\tableofcontents

\section{Background and where we started}

The K=3 CRRA REE operator $\Phi$ takes a conjectured price function
$P(u_1,u_2,u_3)$ and returns the next-iterate. The chebroots/bisect
implementation we built initially floored at $\|F\|_\infty \sim 1.7\!
\times\!10^{-3}$ across all solvers we tried (\texttt{fp\_analysis.pdf}),
diagnosed as structural from root-extraction non-smoothness. The
tabulated-$\mu$ refactoring (\texttt{tab\_explain.pdf}, $\sim 90\times$
speedup) was found to be \emph{not} the cause of the floor: replaying a
legacy operator that hit machine $\varepsilon$ on its own grid and
adding only the lookup table on top still reached $6.4\!\times\!10^{-15}$
under Newton--Krylov.

\section{Method 1: kernel-band co-area integral}

Replace contour root extraction with a Gaussian band:
\[
   A_v(p, u_k) \;=\; \int\!\!\int K_h\!\big(P(u_a, u_b, u_k) - p\big)\,
                       f_v(u_a) f_v(u_b)\, du_a\, du_b,
   \quad K_h(z) = e^{-z^2/(2h^2)}.
\]
With $P$ as a tensor Chebyshev polynomial and the integral evaluated
on a tensor-GL quadrature grid, $A_v$ is an \emph{analytic} function of
the Cheb coefficients of $P$. $\Phi$ is then analytic in $P$, so
Anderson and Newton--Krylov converge quadratically.

\paragraph{Result.} Machine $\varepsilon$ from both cold and warm starts
at $h \in [0.20, 0.50]$, $G_p{=}121$, $N{=}6$ (Anderson floor
$\sim 10^{-16}$, NK confirms to $10^{-15}$). $\sim 7$\,ms per $\Phi$ at
$G_p{=}121$.

\paragraph{Caveat.} The kernel-band FP \emph{depends on $h$}. The joint
$(N{\to}\infty, h{\to}0)$ limit (\texttt{kern\_tab.pdf}) gives slope
$\alpha^* \approx 0.22$ and deficit $1{-}R^2 \approx 0.20$. Going below
$h \approx 0.10$ at $N{=}12$ fails (the band tightens past the grid
resolution and the operator approaches the discontinuous strict-$h{=}0$
limit).

\section{Method 2: partition-of-unity (POU) at strict $h=0$}

NO kernel, NO bandwidth. The strict co-area integral
$A_v = \int_{P=p} f_v f_v / |\nabla P|\, d\sigma$ is made smooth by
representing it as the sum of two parameterizations:
\[
   A_v = A_v^{(a)} + A_v^{(b)},
\]
where $A_v^{(b)}$ parameterizes the contour by $u_a$ (root in $u_b$
at each GL node) and $A_v^{(a)}$ by $u_b$ (root in $u_a$). Each term is
weighted by a partition-of-unity factor
$w_a = d_a^2 / (d_a^2 + d_b^2)$, $w_b = d_b^2 / (d_a^2 + d_b^2)$, with
$w_a + w_b = 1$. As $|d_b| \to 0$ (turning point along $u_b$), $w_b \to 0$
and the $1/|d_b|$ singularity is killed; the contribution shifts to
$A^{(a)}$.

This is a direct port of the K=3 \texttt{hfree\_smooth} legacy operator
(\texttt{projects/REZN/solved\_fixed\_points/k3\_hfree\_smooth/}) to
the Cheb-tab architecture. The atanh-stretched $u$-coordinate makes
$f_v(u(\xi{=}\pm 1)) = 0$ exactly, so roots crossing the box edge
contribute zero amplitude smoothly --- no bandwidth or cutoff needed
for the boundary behaviour.

\subsection{Implementation: scan-and-bisect all-roots in numba}

All-roots-in-$[-1,1]$ via dense sign-change scan
($8N$ subintervals) + 80-iter bisection, fully numba-jit'd:
\texttt{cheby\_numba\_pou\_tab.py}.

Per-$\Phi$ wall time at $G_p{=}121$:
\begin{center}
\begin{tabular}{rrc}
\toprule
$N$ & $G$ & ms / $\Phi$ \\
\midrule
6  & 7  & 52  \\
8  & 9  & 93  \\
10 & 11 & 188 \\
\bottomrule
\end{tabular}
\end{center}

\subsection{Results: correct strict-$h{=}0$ FP, floor $\sim 4\!\times\!10^{-3}$}

\begin{center}
\begin{tabular}{cccc}
\toprule
$N$ ($G$) & best $\|F\|_\infty$ & slope $\alpha^*$ & deficit \\
\midrule
6  (G=7)  & $5.7\!\times\!10^{-3}$ & 0.336 & 0.139 \\
\textbf{8  (G=9)}  & $\mathbf{3.9\!\times\!10^{-3}}$ & \textbf{0.366} & \textbf{0.177} \\
10 (G=11) & $1.3\!\times\!10^{-2}$ & 0.439 & 0.165 (cold-start basin issue) \\
\midrule
Legacy hfree\_smooth (G=9, spline+exact NK) & $4\!\times\!10^{-5}$ & 0.364 & 0.173 \\
\bottomrule
\end{tabular}
\end{center}

\textbf{The POU N=8 slope and deficit agree with the legacy to three
digits} --- the same strict-$h{=}0$ REE. Our generic Anderson+NK
floors $\sim 100\times$ above the legacy because the legacy used
exact-Jacobian Newton steps that contract quadratically.

\section{Investigations into the $4\!\times\!10^{-3}$ POU floor}

\subsection{Pure dense Newton with finite-difference Jacobian}

40-DOF system. FD Jacobian costs $\sim 40\!\cdot\!\Phi = 2$\,s per step.
Sweep $\varepsilon_{\rm FD} \in \{10^{-5}, 10^{-6}, 10^{-7}, 10^{-8}\}$,
warm-start from POU Anderson best:

\begin{center}
\begin{tabular}{cc}
\toprule
$\varepsilon_{\rm FD}$ & best $\|F\|_\infty$ \\
\midrule
$10^{-5}$ & $3.70\!\times\!10^{-3}$ \\
$10^{-6}$ & $\mathbf{2.71\!\times\!10^{-3}}$ \\
$10^{-7}$ & $3.79\!\times\!10^{-3}$ \\
$10^{-8}$ & $3.82\!\times\!10^{-3}$ \\
\bottomrule
\end{tabular}
\end{center}

Marginal improvement over Anderson (4.1e-3). Line-search alpha collapses
to $10^{-9}$, confirming the Jacobian carries non-smoothness noise:
\textbf{the operator, not the solver, is the bottleneck.}

\subsection{Chebroots (companion-matrix all-roots) vs scan-and-bisect}

Hypothesis: my numba scan-and-bisect misses roots near the box
boundary, contributing to the floor. Test: replace with
\texttt{numpy.polynomial.chebyshev.chebroots} (companion matrix
eigenvalues, finds all real roots exactly):

\begin{center}
\begin{tabular}{lcc}
\toprule
all-roots method & ms / $\Phi$ & max $|\Delta P|$ vs scan-bisect \\
\midrule
scan-and-bisect (numba) & 50 & --- (reference) \\
chebroots (numpy) & 2200 & $3.7\!\times\!10^{-4}$ \\
\bottomrule
\end{tabular}
\end{center}

\textbf{The per-$\Phi$ difference is the same order as our floor.}
Strong evidence that scan-and-bisect misses roots near the boundary
and this is what limits POU convergence. The natural fix is to port
the chebroots-based all-roots to numba (companion matrix eigenvalues
are well-supported in numba's \texttt{linalg.eigvals}) or to dramatically
increase the scan resolution near the boundary.

\section{Path forward to machine $\varepsilon$ at strict $h=0$}

If the chebroots experiment confirms much lower POU floor, the next
implementation step is straightforward:
\begin{enumerate}
  \item Port chebroots (companion matrix) to numba — Cheb's
        comrade matrix at $N\!=\!6$ is $6\!\times\!6$, eigenvalue
        extraction is fast and supported in numba.
  \item Re-run POU Anderson+NK; expect floor $\to$ $10^{-7}$ or lower
        if root-finding completeness was the bottleneck.
\end{enumerate}

If FD-Newton still floors above $10^{-7}$ after fixing the root-finding,
the remaining route is \textbf{analytic-Jacobian Newton} (the legacy
path):
\begin{enumerate}
  \item Implement complex-step differentiation: replace all float64 in
        the POU operator with complex128, compute $\partial F / \partial x_j
        = \mathrm{Im}[F(x + i\,h\,e_j)] / h$ with $h \!=\! 10^{-30}$.
        Exact to machine $\varepsilon$ (no FD cancellation).
  \item Or: derive $\partial A_v / \partial P_{\rm coeffs}$ in closed
        form. POU is differentiable (Newton-tracked roots give smooth
        $\partial \xi_b^* / \partial P_{\rm coeffs} = -\partial_{\rm coeff} P /
        \partial_{\xi_b} P$ at the root); the per-$P$-coeff sensitivity
        of $A_v$ chains through the partition weights, Bayes update,
        and CRRA clearing.
\end{enumerate}

The kernel-band Cheb-tab \emph{already} reaches machine $\varepsilon$
for the smoothed operator. The remaining question is whether the
\emph{strict} $h{=}0$ FP (which is the actual REE) is reachable at
machine $\varepsilon$ in pure float64 without arb-precision arithmetic.
The legacy $k3\_hfree\_100$ work shows it can be nailed to 100 decimals
in arb; whether float64 is enough is an empirical question.

\end{document}
